\section{Introduction}
Ultrasound image analysis has become an important pattern recognition problem because ultrasound is widely used in screening, diagnosis support, and follow-up examination for multiple organs. Compared with computed tomography and magnetic resonance imaging, ultrasound is inexpensive, radiation-free, portable, and suitable for repeated acquisition. These advantages make ultrasound attractive for computer-aided diagnosis, but they also make the recognition problem visually difficult. Speckle noise, low contrast, blurred lesion boundaries, operator dependency, and device-specific appearance shifts can make benign--malignant discrimination substantially harder than standard natural-image classification. A useful ultrasound classifier therefore needs to preserve local lesion details, use surrounding tissue context, and remain compact enough for edge-side feasibility studies.

\begin{figure}[!t]
    \centering
    \includegraphics[width=0.95\textwidth]{fig_teaser_tradeoff.png}
    \caption{Accuracy--complexity trade-off on the three public ultrasound datasets. The vertical axis reports the mean AUC over TN5000, BUSI, and AUL; the horizontal axis reports MACs for one forward pass; bubble area denotes model parameters.}
    \label{fig:teaser_tradeoff}
\end{figure}

Public datasets now make it possible to evaluate ultrasound classifiers beyond a single organ. TN5000 provides a large public benchmark for thyroid nodule analysis \cite{Zhang_2025}. BUSI is widely used for breast ultrasound lesion classification \cite{Al_Dhabyani_2020}. The Annotated Ultrasound Liver (AUL) dataset broadens the setting to liver ultrasound \cite{aul_dataset}. These datasets are not interchangeable: they differ in anatomy, acquisition appearance, class prevalence, annotation style, and split protocol. This makes them useful for testing whether an architectural idea is only tuned to one benchmark or remains competitive under different public ultrasound settings.

Recent deep learning methods have improved ultrasound analysis, but three gaps remain important for practical lesion-classification systems. First, many studies report strong results on one dataset while providing limited evidence that the same architecture remains useful under other public ultrasound benchmarks. Second, ROI-based classifiers often rely on annotation-derived crops, which is acceptable for controlled classifier evaluation but leaves the inference-time localization step under-specified. Third, increasingly complex backbones can obscure the practical trade-off between accuracy and computational cost. This trade-off is not a secondary issue in ultrasound: many support scenarios require stable inference under constrained hardware and repeatable preprocessing.

This paper addresses these gaps with UL-TransXNet, a compact ROI framework built from the TransXNet family for ultrasound lesion classification. The method should be understood as a task-specific organization of ROI construction, compact local--global feature modeling, validation-only model-family selection, automatic-ROI workflow evaluation, and mobile prototype measurement. It is not positioned as a new universal medical backbone. TN5000 is treated as the primary large benchmark. BUSI and AUL provide additional dataset-specific in-domain evaluations on breast and liver ultrasound. Structure and module ablations quantify how the TransXNet-family refinements behave under the unified protocol, while automatic-ROI experiments test how much performance is retained when the crop is supplied by a detector rather than by an oracle annotation. A mobile prototype then examines whether the resulting model family can occupy a practical accuracy--latency region.

The main contributions are summarized as follows:
\begin{itemize}[leftmargin=1.5em]
    \item We present UL-TransXNet as a compact ROI classification framework that combines expanded lesion-centered crops with hierarchical local--global feature modeling for ultrasound benign--malignant discrimination.
    \item We evaluate the framework on TN5000, BUSI, and AUL under dataset-specific training, and we explicitly separate multi-dataset in-domain competitiveness from stronger but unsupported claims of cross-organ generalization or uniform superiority.
    \item We audit the final model-family choice with validation-only selection across the three public datasets and provide source-file checksums for the contributing result CSVs, reducing the risk of test-set model selection.
    \item We report automatic-ROI and Android prototype experiments to connect the controlled ROI classifier to localization and deployment workflow evidence, while keeping detector novelty and clinical deployment claims outside the supported conclusion.
\end{itemize}

The remainder of this paper is organized as follows. Section~\ref{sec:related_work} reviews related work. Section~\ref{sec:method} introduces the proposed method. Section~\ref{sec:experimental_setup} presents the datasets and protocol. Section~\ref{sec:results} reports the main experimental findings. Section~\ref{sec:discussion} discusses the implications and limitations, and Section~\ref{sec:conclusion} concludes the paper.

